% GENERATED FILE. Edit canonical lessons, then run npm run generate:books.
% canonical-source-hash: fnv1a64:e62254047de29557
% canonical-lessons: HI-C473-batuni, HI-C473-cup, HI-C473-gap, HI-C473-barbarana, HI-C473-maun

\chapter{Talkers and Quiet Ones}
\label{ch:hi-talkers-and-quiet-ones}
\hlchaptermodality{473}

\begin{chapteropening}
\textbf{By the end of this chapter:} \emph{I can say whether someone talks a lot or keeps quiet.}

Complete the last lesson of chapter 473: I can say whether someone talks a lot or keeps quiet.
\end{chapteropening}

\section[\texorpdfstring{\hi{बातूनी} (bātūnī)}{बातूनी (bātūnī)}]{\hi{बातूनी} (bātūnī) — talkative}

\label{lesson:HI-C473-batuni}



\begin{quote}
Before the new one: say the Hindi for dedicated, then the Hindi for alert.
\end{quote}

\subsection*{You'll want to know: \hi{बातूनी}}
\textbf{\hi{बातूनी}} — \emph{bātūnī} — \textquotedblleft{}talkative\textquotedblright{}.

\begin{quote}
\hi{मेरी} \hi{मौसी} \hi{बहुत} \hi{बातूनी} \hi{हैं}, \hi{घंटों} \hi{बात} \hi{करती} \hi{हैं।}
\end{quote}

\begin{quote}
\emph{merī mausī bahut bātūnī haiṁ, ghaṇṭoṁ bāt kartī haiṁ.}
\end{quote}

\begin{quote}
My aunt is very talkative; she talks for hours.
\end{quote}

\begin{cousinweb}
\textbf{\hi{बातूनी}} (\emph{bātūnī}) is \emph{bāt}, \textquotedblleft{}talk\textquotedblright{}, with the ending \emph{-ūnī}; CDIAL lists Hindi \emph{bāt}, \textquotedblleft{}talk, matter, thing\textquotedblright{}, under Sanskrit \emph{vārtta} (CDIAL 11564).
\end{cousinweb}

\begin{grammarlens}[title={What you've built}]
One more word for describing what you lived and felt.
\end{grammarlens}

\subsection*{Your turn}
\begin{itemize}
\raggedright
  \item \emph{Say it:} \emph{bātūnī}
  \item \emph{Say it:} \emph{bātūnī}, once more
  \item \emph{Say it:} \emph{merī mausī bahut bātūnī haiṁ, ghaṇṭoṁ bāt kartī haiṁ.} — My aunt is very talkative; she talks for hours.
  \item \emph{Recall:} say the Hindi for dedicated, then the Hindi for alert, then say \emph{bātūnī} again
\end{itemize}

\begin{tcolorbox}[breakable,colback=teal!4,colframe=teal!35!black,title={Before you move on}]
What does \hi{बातूनी} mean? (\textquotedblleft{}Talkative\textquotedblright{}.) Say it once more.
\end{tcolorbox}
\section[\texorpdfstring{\hi{चुप} (cup)}{चुप (cup)}]{\hi{चुप} (cup) — silent; quiet}

\label{lesson:HI-C473-cup}



\begin{quote}
Before the new one: say the Hindi for alert, then the Hindi for talkative.
\end{quote}

\subsection*{You'll want to know: \hi{चुप}}
\textbf{\hi{चुप}} — \emph{cup} — \textquotedblleft{}silent; quiet\textquotedblright{}.

\begin{quote}
\hi{माँ} \hi{ने} \hi{डाँटा} \hi{तो} \hi{बच्चा} \hi{चुप} \hi{हो} \hi{गया।}
\end{quote}

\begin{quote}
\emph{māṁ ne ḍā̃ṭā to baccā cup ho gayā.}
\end{quote}

\begin{quote}
When his mother scolded him, the child went quiet.
\end{quote}

\begin{cousinweb}
\textbf{\hi{चुप}} (\emph{cup}) is a tadbhava: CDIAL lists Hindi \emph{cup}, \textquotedblleft{}silent\textquotedblright{}, under \emph{*cuppa}, a form it reconstructs (CDIAL 4864). \emph{cup raho!} is \textquotedblleft{}be quiet!\textquotedblright{}.
\end{cousinweb}

\begin{grammarlens}[title={What you've built}]
One more word for describing what you lived and felt.
\end{grammarlens}

\subsection*{Your turn}
\begin{itemize}
\raggedright
  \item \emph{Say it:} \emph{cup}
  \item \emph{Say it:} \emph{cup}, once more
  \item \emph{Say it:} \emph{māṁ ne ḍā̃ṭā to baccā cup ho gayā.} — When his mother scolded him, the child went quiet.
  \item \emph{Recall:} say the Hindi for alert, then the Hindi for talkative, then say \emph{cup} again
\end{itemize}

\begin{tcolorbox}[breakable,colback=teal!4,colframe=teal!35!black,title={Before you move on}]
What does \hi{चुप} mean? (\textquotedblleft{}Silent; quiet\textquotedblright{}.) Say it once more.
\end{tcolorbox}
\section[\texorpdfstring{\hi{गप} (gap)}{गप (gap)}]{\hi{गप} (gap) — idle chat; gossip}

\label{lesson:HI-C473-gap}



\begin{quote}
Before the new one: say the Hindi for talkative, then the Hindi for silent.
\end{quote}

\subsection*{You'll want to know: \hi{गप}}
\textbf{\hi{गप}} — \emph{gap} — \textquotedblleft{}idle chat; gossip\textquotedblright{}.

\begin{quote}
\hi{शाम} \hi{को} \hi{दादा} \hi{दोस्तों} \hi{के} \hi{साथ} \hi{गप} \hi{मारते} \hi{हैं।}
\end{quote}

\begin{quote}
\emph{śām ko dādā dostoṁ ke sāth gap mārte haiṁ.}
\end{quote}

\begin{quote}
In the evening my grandfather chats with his friends.
\end{quote}

Feminine.

\begin{cousinweb}
\textbf{\hi{गप}} (\emph{gap}) is a tadbhava: CDIAL lists Hindi \emph{gap}, \textquotedblleft{}tittle-tattle\textquotedblright{}, under \emph{*gappa}, a form it reconstructs (CDIAL 4022). \emph{gap mārnā} is to chat idly.
\end{cousinweb}

\begin{grammarlens}[title={What you've built}]
One more word for describing what you lived and felt.
\end{grammarlens}

\subsection*{Your turn}
\begin{itemize}
\raggedright
  \item \emph{Say it:} \emph{gap}
  \item \emph{Say it:} \emph{gap}, once more
  \item \emph{Say it:} \emph{śām ko dādā dostoṁ ke sāth gap mārte haiṁ.} — In the evening my grandfather chats with his friends.
  \item \emph{Recall:} say the Hindi for talkative, then the Hindi for silent, then say \emph{gap} again
\end{itemize}

\begin{tcolorbox}[breakable,colback=teal!4,colframe=teal!35!black,title={Before you move on}]
What does \hi{गप} mean? (\textquotedblleft{}Idle chat; gossip\textquotedblright{}.) Say it once more.
\end{tcolorbox}
\section[\texorpdfstring{\hi{बड़बड़ाना} (baṛbaṛānā)}{बड़बड़ाना (baṛbaṛānā)}]{\hi{बड़बड़ाना} (baṛbaṛānā) — to grumble; to mutter}

\label{lesson:HI-C473-barbarana}



\begin{quote}
Before the new one: say the Hindi for silent, then the Hindi for idle chat.
\end{quote}

\subsection*{You'll want to know: \hi{बड़बड़ाना}}
\textbf{\hi{बड़बड़ाना}} — \emph{baṛbaṛānā} — \textquotedblleft{}to grumble; to mutter\textquotedblright{}.

\begin{quote}
\hi{काम} \hi{ज़्यादा} \hi{था}, \hi{इसलिए} \hi{नौकर} \hi{बड़बड़ा} \hi{रहा} \hi{था।}
\end{quote}

\begin{quote}
\emph{kām zyādā thā, isliye naukar baṛbaṛā rahā thā.}
\end{quote}

\begin{quote}
There was too much work, so the servant was grumbling.
\end{quote}

Past (she …): \hi{वह} \hi{कुछ} \hi{बड़बड़ाई}, \emph{vah kuch baṛbaṛāī}.

\begin{cousinweb}
\textbf{\hi{बड़बड़ाना}} (\emph{baṛbaṛānā}) is a tadbhava: CDIAL lists Hindi \emph{baṛbaṛānā}, \textquotedblleft{}to grumble\textquotedblright{}, under \emph{*baḍabaḍa}, a form it reconstructs (CDIAL 9122); the doubled syllable is the sound of muttering.
\end{cousinweb}

\begin{grammarlens}[title={What you've built}]
One more word for describing what you lived and felt.
\end{grammarlens}

\subsection*{Your turn}
\begin{itemize}
\raggedright
  \item \emph{Say it:} \emph{baṛbaṛānā}
  \item \emph{Say it:} \emph{baṛbaṛānā}, once more
  \item \emph{Say it:} \emph{kām zyādā thā, isliye naukar baṛbaṛā rahā thā.} — There was too much work, so the servant was grumbling.
  \item \emph{Recall:} say the Hindi for silent, then the Hindi for idle chat, then say \emph{baṛbaṛānā} again
\end{itemize}

\begin{tcolorbox}[breakable,colback=teal!4,colframe=teal!35!black,title={Before you move on}]
What does \hi{बड़बड़ाना} mean? (\textquotedblleft{}To grumble; to mutter\textquotedblright{}.) Say it once more.
\end{tcolorbox}
\section[\texorpdfstring{\hi{मौन} (maun)}{मौन (maun)}]{\hi{मौन} (maun) — silence}

\label{lesson:HI-C473-maun}



\begin{quote}
Before the new one: say the Hindi for idle chat, then the Hindi for to grumble.
\end{quote}

\subsection*{You'll want to know: \hi{मौन}}
\textbf{\hi{मौन}} — \emph{maun} — \textquotedblleft{}silence\textquotedblright{}.

\begin{quote}
\hi{स्कूल} \hi{में} \hi{बच्चों} \hi{ने} \hi{दो} \hi{मिनट} \hi{का} \hi{मौन} \hi{रखा।}
\end{quote}

\begin{quote}
\emph{skūl meṁ baccoṁ ne do minaṭ kā maun rakhā.}
\end{quote}

\begin{quote}
At school the children kept two minutes' silence.
\end{quote}

Masculine.

\begin{cousinweb}
\textbf{\hi{मौन}} (\emph{maun}) is a tatsama, borrowed from Sanskrit \emph{mauna}, \textquotedblleft{}silence, taciturnity\textquotedblright{}, from \emph{muni}, \textquotedblleft{}a silent sage\textquotedblright{} (Monier-Williams). It is the formal, chosen silence; \emph{cup} is the everyday kind.
\end{cousinweb}

\begin{grammarlens}[title={What you've built}]
One more word for describing what you lived and felt.
\end{grammarlens}

\subsection*{Your turn}
\begin{itemize}
\raggedright
  \item \emph{Say it:} \emph{maun}
  \item \emph{Say it:} \emph{maun}, once more
  \item \emph{Say it:} \emph{skūl meṁ baccoṁ ne do minaṭ kā maun rakhā.} — At school the children kept two minutes' silence.
  \item \emph{Recall:} say the Hindi for idle chat, then the Hindi for to grumble, then say \emph{maun} again
\end{itemize}

\begin{tcolorbox}[breakable,colback=teal!4,colframe=teal!35!black,title={Before you move on}]
What does \hi{मौन} mean? (\textquotedblleft{}Silence\textquotedblright{}.) Say it once more.
\end{tcolorbox}
